\subsection{Notation} \label{notation}

We represent a security parameter as $\lambda \in \mathbb{Z}$.
We denote the assignment of an element $y$ to the value $x$ as $y \gets x$,
and sampling an element from some set $S$ uniformly at random as $x \rgets S$.
For a randomized algorithm $\mathit{A}$, we write $x \rgets \mathit{A}()$ to indicate the random variable $x$ that is output from the execution of $\mathit{A}$.
We say that $f$ is a negligible function of $\lambda$ if it holds that
$\lim_{\lambda \to \infty} f(\lambda) \lambda^c = 0$,
for all $ c > 0$~\cite{BonehS2020}.

Let $\GG$ be a cyclic group of prime order $q$, and $\Zq$ be the field of integers modulo $q$.
Let $g$ be a generator of $\GG$.

We use $\set{n}$ to represent the set $\{1,\ldots,n \}$,
and represent integer vectors in bold-face font.
For a set $S$,
we denote ${S \choose t}$ to mean the set
that consists of all size-$t$ subsets of $S$.

\paragraph{Polynomial Interpolation.}
A polynomial $f(x) = a_0 + a_1 x + a_2 x^2 + \ldots + a_{\tcor} x^{\tcor}$
of degree $\tcor$ over a field $\F$ can be interpolated by $\thresh$ points.
Let $\eta$ be a set of $\thresh$ distinct integers
corresponding to the $x$-coordinates $x_i \in \F, i \in \eta$, of these points.
Then the Lagrange polynomial $L_i(x)$ has the form:

\begin{equation}\label{eqn:lagrange}
L_i(x) = \prod_{j \in \eta; j \neq i } \frac{x-x_j}{x_i - x_j}
\end{equation}

Given a set of $\thresh$ points $(x_i, f(x_i))_{i \in [\eta]}$, any point $f(x_\ell)$ on the polynomial $f$ can be determined by Lagrange interpolation as follows:

\[ f(x_\ell) = \sum_{k \in \eta} f(x_k) \cdot L_k(x_\ell) \]

\paragraph{Game-Based Security.}
An experiment (or game) $\textsf{Exp}_{\mathcal{C}, \adv}^{\mathsf{sec}}$ is
played with respect to a security notion $\mathsf{sec}$,
an adversary $\adv$ modeled by a probabilistic (randomized) algorithm,
and a cryptographic primitive $\mathcal{C}$.
The experiment is initialized within a main procedure,
whose output is the output of the game.
The goal of $\adv$ is to win the game,
the conditions under which is defined within the experiment.
The probability of $\adv$'s success thus defines how secure $\mathcal{C}$ is with respect to the security notion $\mathsf{sec}$.

\paragraph{PPT Adversaries.}
We say that an adversary $\adv$ is probabilistic polynomial-time (PPT) if it is run on
random coins and its execution is performed in time that is polynomial relative to the size of its input.
In other words,
there exists a polynomial $f$,
such that for every input $x \in \zo^n$,
$\adv(x)$ terminates within at most $f(n)$ steps.

\paragraph{Group Generation.}
Let $\GroupGen$ be a probabilistic polynomial-time algorithm that takes as input a security parameter
$1^\usecp$ and outputs a group description $(\Gr, q, g)$ consisting of a group $\Gr$ of order $q$,
where $q$ is a $2^\secp$-bit prime,
and a generator $g$ of $\Gr$.
